% !TeX root = writeup.tex 


\section{Introduction}

Machine learning commonly considers static objectives defined on a snapshot of
the population at one instant in time; consequential decisions, in contrast,
reshape the population over time. Lending practices, for example, can shift the
distribution of debt and wealth in the population. Job advertisements 
allocate opportunity. School admissions shape the level of education in
a community.

Existing scholarship on fairness in automated decision-making criticizes
unconstrained machine learning for its potential to \emph{harm} historically
underrepresented or disadvantaged groups in the
population~\citep{whitehouse16,barocasselbst16}. Consequently, a variety of
\emph{fairness criteria} have been proposed as constraints on standard learning
objectives. Even though, in each case, these constraints are clearly intended to
\emph{protect} the disadvantaged group by an appeal to intuition, a rigorous
argument to that effect is often lacking.

In this work, we formally examine under what circumstances fairness criteria do
indeed promote the long-term well-being of disadvantaged groups measured in
terms of a temporal variable of interest. Going beyond the standard
classification setting, we introduce a one-step feedback model of decision-making
that exposes how decisions change the underlying population over time.

Our running example is a hypothetical lending scenario. There are two groups in
the population with features described by a summary statistic, such as a
\emph{credit score}, whose distribution differs between the two groups. The bank can
choose thresholds for each group at which loans are offered. While
group-dependent thresholds may face legal challenges \citep{ross2006}, they are
generally inevitable for some of the criteria we examine. The impact of a
lending decision has multiple facets. A default event not only diminishes profit for
the bank, it also worsens the financial situation of the borrower as reflected
in a subsequent decline in credit score. A successful lending outcome leads to
profit for the bank and also to an increase in credit score for the borrower.

When thinking of one of the two groups as disadvantaged, it makes sense to ask
what lending policies (choices of thresholds) lead to an expected improvement in
the score distribution within that group. An unconstrained bank would maximize
profit, choosing thresholds that meet a break-even point above which it is
profitable to give out loans. One frequently proposed fairness criterion,
sometimes called demographic parity, requires the bank to lend to both groups at
an equal rate. Subject to this requirement the bank would continue to maximize
profit to the extent possible. 
Another criterion, originally called equality of opportunity, equalizes
the \emph{true positive rates} between the two groups, thus requiring the bank to lend in
both groups at an equal rate among individuals who repay their loan. Other
criteria are natural, but for clarity we restrict our attention to these three.

Do these fairness criteria benefit the disadvantaged group? When do they show a
clear advantage over unconstrained classification? Under what circumstances does
profit maximization work in the interest of the individual? These are important
questions that we begin to address in this work.

\subsection{Contributions}
We introduce a one-step feedback model that allows us to quantify the long-term
impact of classification on different groups in the population. We represent
each of the two groups~$\popa$ and~$\popb$ by a \emph{score}
distribution~$\dist\supa$ and~$\dist\supb,$ respectively. The support of these
distributions is a finite set~$\X$ corresponding to the possible values that the
score can assume. We think of the score as highlighting one variable of interest
in a specific domain such that higher score values correspond to a
higher probability of a positive outcome.  
An \emph{institution} chooses selection policies~$\pola,
\polb\colon\X\to[0,1]$ that assign to each value in~$\X$ a number representing
the rate of selection for that value. In our example, these policies specify
the lending rate at a given credit score within a given group. The institution
will always maximize their utility (defined formally later) subject to either
(a) no constraint, or (b) equality of selection rates, or (c) equality of true
positive rates.

We assume the availability of a function~$\chg\colon\X\to\R$ such that $\chg(x)$ provides the
expected change in score for a selected individual at score $x$.  The
central quantity we study is the expected difference in the mean
score in group~$j\in\{\popa, \popb\}$ that results from an institutions policy, $\delmean\supj$ defined formally in Equation~\eqref{eqn:delmean}. 
When modeling the problem, the expected mean difference can also absorb external
factors such as ``reversion to the mean'' so long as they are mean-preserving.
Qualitatively, we distinguish between \emph{long-term improvement}
($\delmean\supj>0$), \emph{stagnation} ($\delmean\supj=0$), and \emph{decline}
($\delmean\supj<0$).
Our findings can be summarized as follows:

\begin{enumerate} 
	\item Both fairness criteria (equal selection rates, equal true
	positive rates) can lead to all possible outcomes (improvement, stagnation, and decline) in natural parameter regimes.  We provide a complete characterization of when each criterion
	leads to each outcome in Section \ref{sec:results}.
	\begin{itemize}
		\item There are a class of settings where equal
		selection rates cause decline, whereas equal true positive rates do not (Corollary \ref{cor:avoid_harm}),
		\item Under a mild assumption, the institution's optimal unconstrained selection policy can never lead to decline (Proposition \ref{prop:mp_noharm}). 
	\end{itemize}

\item We introduce the notion of an \emph{outcome curve} (Figure~\ref{fig:outcome_curve}) which succinctly describes the different regimes in which one criterion is preferable over the others. 

\item We perform experiments on FICO credit score data from 2003 and show that under various models of bank utility and score change, the outcomes of applying fairness criteria are in line with our theoretical predictions.

\item
We discuss how certain types of measurement error (e.g., the bank
underestimating the repayment ability of the disadvantaged group) affect our
comparison. We find that measurement error narrows the regime in which fairness
criteria cause decline, suggesting that measurement should be a factor when
motivating these criteria.
\item We consider alternatives to hard fairness constraints.
\begin{itemize}
	\item We evaluate the optimization problem where fairness criterion is a regularization term in the objective. Qualitatively, this leads to the
	same findings.
	\item We discuss the possibility of optimizing for group score improvement $\delmean\supj$ directly subject to institution utility constraints. The
	resulting solution provides an interesting possible alternative to existing fairness criteria.
\end{itemize}

\end{enumerate}

We focus on the impact of a selection policy over a single epoch. The
motivation is that the designer of a system usually has an understanding of the
time horizon after which the system is evaluated and possibly redesigned.
Formally, nothing prevents us from repeatedly applying our model and tracing
changes over multiple epochs. In reality, however, it is plausible that over greater time
periods, economic background variables might dominate the effect of selection.


Reflecting on our findings, we argue that careful temporal modeling is necessary
in order to accurately evaluate the impact of different fairness criteria on the
population. Moreover, an understanding of measurement error is important in
assessing the advantages of fairness criteria relative to unconstrained
selection. Finally, the nuances of our characterization underline how
intuition may be a poor guide in judging the long-term impact of fairness
constraints.

\subsection{Related work}

%\lydia{Add references}

Recent work by \citet{hu18shortterm} considers a model for long-term outcomes and fairness in the labor market. They propose imposing the demographic parity constraint in a \emph{temporary} labor market in order to provably achieve an equitable long-term equilibrium in the \emph{permanent} labor market, reminiscent of economic arguments for affirmative action~\citep{foster1992economic}. The equilibrium analysis of the labor market dynamics model allows for specific conclusions relating fairness criteria to long term outcomes. Our general framework is complementary to this type of domain specific approach.

\citet{fuster2017predictably}
consider the problem of fairness in credit markets from a different perspective. Their goal is to study the effect of machine learning on interest rates in different groups at an equilibrium, under a static model without feedback.

\citet{ensign2017runaway} consider feedback loops in predictive
policing, where the police more heavily monitor high crime neighborhoods,
thus further increasing the measured number of crimes in those neighborhoods.
While the work addresses an important temporal phenomenon using the theory of urns, it is rather different from our one-step feedback model both conceptually and technically.

Demographic parity and its related formulations have been considered in numerous
papers~\citep[e.g.][]{Calders2009,zafar17a}. \citet{hardt16equality} introduced the equality of
opportunity constraint that we consider and demonstrated limitations of a broad
class of criteria. \citet{kleinberg16inherent} and~\citet{chould16fair} point out
the tension between ``calibration by group'' and equal true/false positive
rates. These trade-offs carry over to some extent to the case where we only
equalize true positive rates~\citep{weinberger2017calibration}.

A growing literature on fairness in the ``bandits'' setting of learning~\citep[see][\emph{et sequelae}]{roth2016bandits} deals with online decision making that ought not
to be confused with our one-step feedback setting. Finally, there has been much work in the social sciences on analyzing the effect
of affirmative action \citep[see e.g.,][]{keith1985,Kaley2006}.

\subsection{Discussion}

%We now discuss the implications of an outcome-based view towards fairness in machine learning that this paper pursues. In particular, we center our discussion around the following two questions.

In this paper, we advocate for a view toward long-term outcomes in the
discussion of ``fair'' machine learning. 
%
We argue that without a careful model of delayed outcomes, we cannot foresee the
impact a fairness criterion would have if enforced as a constraint on a
classification system.  However, if such an accurate outcome model is available,
we show that there are more direct ways to optimize for positive outcomes than
via existing fairness criteria. We outline such an outcome-based solution 
in Section \ref{sec:outcome_based}. Specifically, in the credit setting, the
outcome-based solution corresponds to giving out more loans to the protected
group in a way that reduces profit for the bank compared to unconstrained profit
maximization, but avoids loaning to those who are unlikely to benefit, resulting
in a maximally improved group average credit score. The extent to which such a
solution could form the basis of successful regulation depends on the accuracy
of the available outcome model.

This raises the question if our model of outcomes is rich enough to faithfully
capture realistic phenomena.  By focusing on the impact that selection has on
individuals at a given score, we model the effects for those \emph{not} selected
as zero-mean. For example, not getting a loan in our model has no negative
effect on the credit score of an individual.\footnote{In reality, a denied
credit inquiry may lower one's credit score, but the effect is small compared to
a default event.} This does not mean that wrongful rejection (i.e., a false
negative) has no visible manifestation in our model. If a classifier has a
higher false negative rate in one group than in another, we expect the
classifier to increase the disparity between the two groups (under natural
assumptions). In other words, in our outcome-based model, the harm of denied
opportunity manifests as growing disparity between the groups.  The cost of a
false negative could also be incorporated directly into the outcome-based model by a simple modification (see
Footnote~\ref{footnote:negative_dec_outcomes}). This may be fitting in some
applications where the immediate impact of a false negative to the individual is
not zero-mean, but significantly reduces their future success probability.

In essence, the formalism we propose requires us to understand the two-variable
causal mechanism that translates decisions to outcomes.  This can be seen as
relaxing the requirements compared with recent work 
on avoiding discrimination through causal reasoning that often required stronger
assumptions~\citep{Kusner17,Nabi17,Kilbertus17}. In particular, these works required knowledge of how
sensitive attributes (such as gender, race, or proxies thereof) causally relate
to various other variables in the data. Our model avoids the delicate modeling
step involving the sensitive attribute, and instead focuses on an arguably more
tangible economic mechanism. Nonetheless, depending on the application, such an
understanding might necessitate greater domain knowledge and additional research
into the specifics of the application. This is consistent with much scholarship
that points to the context-sensitive nature of fairness in machine learning. 

